\section*{NeurIPS Paper Checklist}

\begin{enumerate}

\item {\bf Claims}
    \item[] Question: Do the main claims made in the abstract and introduction accurately reflect the paper's contributions and scope?
    \item[] Answer: \answerYes{}
    \item[] Justification: The abstract and introduction state three contributions—a 470-item benchmark, a multi-method mechanistic analysis (probes, SAEs, attention entropy, activation steering), and evidence for a graded deliberation-intensity dimension—all of which are supported by the experimental results presented in the paper.
\item {\bf Limitations}
    \item[] Question: Does the paper discuss the limitations of the work performed by the authors?
    \item[] Answer: \answerYes{}
    \item[] Justification: The Discussion section explicitly addresses limitations: experiments cover models up to 32B parameters (not frontier-scale), linear probes may miss nonlinear representational structure, and the benchmark is limited to English-language cognitive-bias tasks. We also discuss the inherent limitations of the dual-process operationalization.
\item {\bf Theory assumptions and proofs}
    \item[] Question: For each theoretical result, does the paper provide the full set of assumptions and a complete (and correct) proof?
    \item[] Answer: \answerNA{}
    \item[] Justification: This is an empirical study. The paper does not present theoretical results or proofs.
    \item {\bf Experimental result reproducibility}
    \item[] Question: Does the paper fully disclose all the information needed to reproduce the main experimental results of the paper to the extent that it affects the main claims and/or conclusions of the paper (regardless of whether the code and data are provided or not)?
    \item[] Answer: \answerYes{}
    \item[] Justification: We report all hyperparameters, cross-validation procedures (5-fold stratified CV), statistical methods (bootstrap CIs, BH-FDR correction), control tasks (Hewitt \& Liang random-label baselines, random-direction controls), model identifiers, and layer selections. The appendix contains full experimental details sufficient for independent replication.
\item {\bf Open access to data and code}
    \item[] Question: Does the paper provide open access to the data and code, with sufficient instructions to faithfully reproduce the main experimental results, as described in supplemental material?
    \item[] Answer: \answerYes{}
    \item[] Justification: Code and benchmark will be provided via anonymous link at review time. The 470-item benchmark (JSONL) and all analysis code will be released publicly upon acceptance.
\item {\bf Experimental setting/details}
    \item[] Question: Does the paper specify all the training and test details (e.g., data splits, hyperparameters, how they were chosen, type of optimizer) necessary to understand the results?
    \item[] Answer: \answerYes{}
    \item[] Justification: Section~4 and the appendix specify data splits (5-fold stratified CV by task category and label), probe hyperparameters (logistic regression with L2 regularization, selected via nested CV), SAE reconstruction fidelity thresholds, attention entropy normalization, and activation steering parameters. All models are identified by their HuggingFace IDs and accessed via publicly available weights.
\item {\bf Experiment statistical significance}
    \item[] Question: Does the paper report error bars suitably and correctly defined or other appropriate information about the statistical significance of the experiments?
    \item[] Answer: \answerYes{}
    \item[] Justification: All probe accuracy and AUC metrics are reported with bootstrap 95\% confidence intervals (1000 resamples). Multiple comparisons across layers and models are corrected with Benjamini-Hochberg FDR at $q=0.05$. Multi-seed robustness analyses report mean $\pm$ standard deviation across seeds. The source of variability (CV fold, bootstrap resample, or random seed) is stated for each metric.
\item {\bf Experiments compute resources}
    \item[] Question: For each experiment, does the paper provide sufficient information on the computer resources (type of compute workers, memory, time of execution) needed to reproduce the experiments?
    \item[] Answer: \answerYes{}
    \item[] Justification: The paper reports that all experiments were conducted on NVIDIA B200 GPUs via RunPod, with a total compute budget of approximately 100 GPU-hours. Per-workstream compute breakdowns are provided in the appendix.
\item {\bf Code of ethics}
    \item[] Question: Does the research conducted in the paper conform, in every respect, with the NeurIPS Code of Ethics \url{https://neurips.cc/public/EthicsGuidelines}?
    \item[] Answer: \answerYes{}
    \item[] Justification: The research uses only publicly available open-source models and does not involve human subjects, deception, or personally identifiable information. All work conforms to the NeurIPS Code of Ethics.
\item {\bf Broader impacts}
    \item[] Question: Does the paper discuss both potential positive societal impacts and negative societal impacts of the work performed?
    \item[] Answer: \answerYes{}
    \item[] Justification: The Discussion section addresses broader impacts: positively, our methods could inform runtime monitoring of LLM reasoning fidelity for AI safety; negatively, better understanding of internal deliberation could theoretically aid adversarial manipulation of model reasoning. We assess the risk as low given that the work is interpretability research on open-source models.
\item {\bf Safeguards}
    \item[] Question: Does the paper describe safeguards that have been put in place for responsible release of data or models that have a high risk for misuse (e.g., pre-trained language models, image generators, or scraped datasets)?
    \item[] Answer: \answerNA{}
    \item[] Justification: The paper releases a cognitive-bias benchmark and analysis code, neither of which poses a meaningful risk for misuse. No new models or scraped datasets are released.
\item {\bf Licenses for existing assets}
    \item[] Question: Are the creators or original owners of assets (e.g., code, data, models), used in the paper, properly credited and are the license and terms of use explicitly mentioned and properly respected?
    \item[] Answer: \answerYes{}
    \item[] Justification: All eight models used are open-source and cited with their original publications. Model licenses (Llama 3.1 Community License, Apache 2.0 for Qwen and OLMo, Gemma Terms of Use) are listed in the appendix. Pre-trained SAEs (Llama Scope, Gemma Scope) are cited and their licenses noted.
\item {\bf New assets}
    \item[] Question: Are new assets introduced in the paper well documented and is the documentation provided alongside the assets?
    \item[] Answer: \answerYes{}
    \item[] Justification: The 470-item S1/S2 benchmark is documented with a datasheet describing item construction methodology, task categories, conflict/no-conflict structure, and intended use. The JSONL schema and analysis code are documented in the supplemental material.
\item {\bf Crowdsourcing and research with human subjects}
    \item[] Question: For crowdsourcing experiments and research with human subjects, does the paper include the full text of instructions given to participants and screenshots, if applicable, as well as details about compensation (if any)?
    \item[] Answer: \answerNA{}
    \item[] Justification: This study does not involve crowdsourcing or human subjects. All experiments are conducted on language model activations.
\item {\bf Institutional review board (IRB) approvals or equivalent for research with human subjects}
    \item[] Question: Does the paper describe potential risks incurred by study participants, whether such risks were disclosed to the subjects, and whether Institutional Review Board (IRB) approvals (or an equivalent approval/review based on the requirements of your country or institution) were obtained?
    \item[] Answer: \answerNA{}
    \item[] Justification: No human subjects were involved in this research. IRB approval is not applicable.
\item {\bf Declaration of LLM usage}
    \item[] Question: Does the paper describe the usage of LLMs if it is an important, original, or non-standard component of the core methods in this research? Note that if the LLM is used only for writing, editing, or formatting purposes and does \emph{not} impact the core methodology, scientific rigor, or originality of the research, declaration is not required.
    %this research?
    \item[] Answer: \answerYes{}
    \item[] Justification: LLMs are the subject of study (we analyze internal representations of 8 open-source LLMs). Additionally, we disclose that AI coding assistants (Claude Code) were used for implementation assistance. Neither use impacts the scientific methodology or originality of the research, but we declare both for transparency.
\end{enumerate}
